\section{Appendix: Scope and Qualifications}
\label{sec:scope}

{\small
\textbf{Edition and formulary.} The study treats the Mass of the Twenty-sixth
Sunday in Ordinary Time, Year~A, as the \work{Roman Missal}, Third Edition, and
the \work{Lectionary for Mass} give it for the dioceses of the United States,
in the national calendar alone, for the occurrence of 27~September 2026.
Identity, occurrence and branches are set out in the preceding appendix and
recorded in \texttt{instance/manifest.md} and \texttt{research/context.md}.

\textbf{Text control.} The Missal formulary was read in the Latin of the
\work{Missale Romanum}, editio typica tertia (2002), in a digitally typeset
reproduction, and compared with the list of changes for the 2008 reprint in
\work{Notitiae} 44 (2008), which lists none for this week, and with the
\work{Antiphonary} excerpted from the English Missal (2010). No page of the 2008
reprint and no United States altar book was inspected, so the approved English
orations were not read in any witness, and no chant book was collated. Which
Latin text the Entrance Antiphon's wording follows is not established. The
Lectionary's boundaries rest on three witnesses that agree in every figure: the
national calendar for 2026, the bishops' conference's dated page of readings,
and no.~136 of the Latin \work{Ordo lectionum Missae} (1981); the printed United
States Lectionary was not inspected. Three questions about the text remain
open: whether the United States strophe includes the last clause of Psalm 25:5;
whether the United States Lectionary admits another Alleluia verse from a
common set on these Sundays; and the manuscript evidence for the order of the
sons and the answer of Matthew 21:29--31 beyond what the SBL apparatus reports.
No manuscript was examined.

\textbf{English.} All Scripture is printed in the Douay--Rheims version as
revised by Bishop Challoner, a public-domain study translation of the
Clementine Vulgate. It is not the text proclaimed in the celebration; the
controlling English is that of the \work{Lectionary for Mass for Use in the
Dioceses of the United States of America}, and for the antiphons and orations
that of the \work{Roman Missal}, Third Edition, neither of which is reproduced.
The orations and the two antiphons that adapt Scripture are identified by Latin
incipit and described; the descriptions are not translations and may not be
quoted as the texts' words. Where the Lectionary appoints part of a verse, or
may do so, the rest is marked beside it. English glosses of Latin passages
are the editor's renderings of the Latin printed beside them.

\textbf{Reception.} Each appointed Scripture passage has at least one ancient
and one medieval or later witness read at an exact locus; the witnesses are
those named in the References, and the search behind them is bounded and
recorded in \texttt{research/scope.md}. For the first reading and the second
Communion antiphon the later witness is Aquinas's doctrinal use of the verse,
and for the Entrance Antiphon it is a liturgical commentator on the chant;
three passages rest on one ancient witness each (the psalm on Augustine, the
Entrance Antiphon's verses on Jerome, the second Communion antiphon on
Augustine). Chrysostom, Augustine's sermons, his tractates on John and on the
First Epistle of John, and his expositions of Psalms 17 and 24 were read in the
English of the Nicene and Post-Nicene Fathers; Aquinas's \work{Summa} III in the
English Dominican translation; Bellarmine in O'Sullivan's abridged English of
1866, the Latin not compared; Schuster in a text layer of the 1927 English.
Jerome on Matthew, Ezekiel and Daniel, Hilary on Matthew and on Psalm 118,
Augustine on Psalm 118, and Aquinas on Matthew, John and Philippians and in
\work{Summa} I, q.~25 were read on page images of the printings named; Gregory
and \work{Summa} II-II, q.~30 in electronic Latin texts. No witness was collated
with a modern critical edition. Not opened, among others: Origen and Rabanus
Maurus on Matthew; Theodoret, Albert and Hugh of Saint-Cher on Ezekiel;
Cassiodorus and Theodoret on the psalms; Ambrose on Psalm 118; Chrysostom and
Cyril of Alexandria on John; Bede on the First Epistle of John; Rupert of Deutz
and Berno of Reichenau on the Entrance Antiphon; and the commentary of
\work{The Liturgical Year} on the Collect. Nothing is said here about who first
proposed any reading, about any reading being the Fathers' common view, or
about the order in time or the mutual dependence of the witnesses. No
exposition of the parable of the two sons by Augustine was found in the texts
of his searched, the English sermons and tractates and the Latin
\work{Quaestiones evangeliorum}; the search was literal and does not reach his
works not held. No commentary on the Prayer over the Offerings or the Prayer
after Communion was found, and a search of the sacramentaries held found no
earlier witness to the first prayer or to the present ending of the second,
which does not show that none exists. Schuster's comments on the Collect and
the Entrance Antiphon belong to Masses that differ from this formulary in their
other elements, and he is cited only for those two texts.

\textbf{Relations among the elements.} Three relations are stated by the
liturgical books: the first reading is chosen in relation to the Gospel
(\work{Praenotanda} 106), the psalm responds to the first reading
(\work{General Instruction} 61), and the acclamation greets the Gospel
(\work{General Instruction} 62). The second reading and the Gospel each belong
to a semi-continuous course, and the Missal's antiphons and orations serve all
three years. The recurrence of the words for way, for remembering and for
humility across the texts is an observation about the texts. One witness cited
joins two of the Mass's texts himself: Augustine cites Ezekiel 18:22 in
expounding Psalm 118:49. Every other connection drawn between one element and
another is the editor's synthesis, grounded in a witness's statement about his
own passage where one is cited and otherwise an editorial proposal resting on
the wording of the texts. No witness cited comments on this Mass.

\textbf{The Council's declaration.} \work{Nostra aetate} 4 was read in the
English of the Holy See's website; its Latin was not compared. It is cited as
the Church's teaching on how the identification of the two sons with two
peoples is presented, not as a judgment on the exegesis of any Father.

\textbf{Chronology.} The Date cells print the shared chronology corpus's
answers with their relations and alternatives unchanged. The first reading's
answers date the captivity and Ezekiel's ministry, not the oracle of chapter
18, and the A.M.\ figure is Haydock's report of Usher, whose page was not read
here. The Entrance Antiphon's composition answer is one author's judgment on
the Greek additions to Daniel, and his own article reports the contrary view of
most Catholic writers of his day. The Psalter's common composition bound is
not a date for either psalm; for Psalm 25 it was read in the current public
text of the introduction it cites, not in the registered state. The Gospel's
and the Johannine answers are composition dates of the books. Writing sites,
first readers and narrative settings are kept apart; the source checks and
their limits are recorded in \texttt{research/scope.md}.

\textbf{Rights.} The Douay--Rheims, the Clementine Vulgate, the Nicene and
Post-Nicene Fathers, O'Sullivan's Bellarmine, the English Dominican
\work{Summa}, the volumes of Migne, the Vienna Hilary, the Venice and Li\`ege
printings of Aquinas, Schuster's English, Wilson's and Feltoe's editions and
the \work{Catholic Encyclopedia} are in the public domain. Gregory's Latin is
quoted in short phrases from a licensed electronic edition; the Leonine
\work{Summa}, whose rights are unresolved, and the Corpus Thomisticum text of
II-II, q.~30 are quoted in short passages.
The Latin and English of the Missal, the \work{Ordo lectionum Missae}, the
Lectionary, the \work{General Instruction}, the national calendar, the
\work{Nova Vulgata}, \work{Nostra aetate} in English and the introduction to the
Psalms of the New American Bible are protected; they are cited, and only
incipits, the three titles of the \work{Ordo}, short phrases of the Council's
declaration and a few words needed by the argument are quoted.

\textbf{Records.} The composition of the target is audited in
\texttt{propers/verified.md}; the reception search, loci, disagreements and
negative results in \texttt{research/scope.md}; and the reasoning behind the
three interpretations, with the class of every cross-element link, in
\texttt{research/interpretations.md}. The research records received an
independent review before the study was written.
\par}

\section*{References}
% A starred heading sets no mark, so without this the Appendix above would go
% on heading the pages of the bibliography. The heading stays starred because
% check-content-preflight finds the References section by that exact form.
\runninghead{References}
\addcontentsline{toc}{section}{References}
\label{sec:references}

{\small
\subsection*{Liturgical books, Bibles and texts}
\begin{itemize}[itemsep=0pt,topsep=.25em]
\item \work{Missale Romanum}, editio typica tertia (Typis Vaticanis, 2002), \latin{Dominica XXVI ``per annum''}; the rubrics of \latin{Tempus per annum} 1--6; the \latin{Normae universales de anno liturgico et de calendario} 59--60 and the \latin{Calendarium Romanum generale} at 27~September, as printed in that Missal.
\item \work{Variationes et additiones in Missali Romano}, \work{Notitiae} 44 (2008), pp.~371--372.
\item \work{The Roman Missal}, Third Edition, for Use in the Dioceses of the United States of America (2011), as represented by the \work{Antiphonary: Excerpted from the Roman Missal} (International Commission on English in the Liturgy, 2010), p.~79.
\item \work{Ordo lectionum Missae}, editio typica altera (Libreria Editrice Vaticana, 1981): \work{Praenotanda} 80, 105--107; \latin{Tabella II}; no.~136, \latin{Dominica vigesima sexta}, Year~A.
\item \work{Lectionary for Mass for Use in the Dioceses of the United States of America}, second typical edition, vol.~I, no.~136, as witnessed by the United States Conference of Catholic Bishops, \work{Liturgical Calendar for the Dioceses of the United States of America} 2026, pp.~5 and 39, and its page of readings for 27~September 2026, \url{https://bible.usccb.org/bible/readings/092726.cfm}.
\item \work{General Instruction of the Roman Missal}, United States edition of 2011 with the emendations of 2021, nos.~48, 51, 53, 61--64, 68, 74, 87, 363--365.
\item The Holy Bible, Douay--Rheims version, revised by Bishop Richard Challoner (Project Gutenberg text): Exod 19:8; 24:3, 7; 4~Kgs 24:8--12; Ps 24; 94:8; 118:49--56; Eccles 5:4; Jer 25:11--12; 31:19; Ezek 1:1--3; 8:1; 18; 20:1; 29:17; Dan 3:1--45; Mt 21:23--46; Jn 10:11--30; Rom 2:13; 8:17; 1~Cor 11:26; Phil 1:1--2:15; 1~Pet 2:21; 2~Pet 2:21; 1~Jn 3:11--18.
\item \work{Biblia Sacra Vulgatae editionis} (Clementine Vulgate), eBible.org text: Ps 24:4--10; 118:49--50; Ezek 18:22, 25--29; Dan 3:27--43; Mt 21:28--32; Jn 10:26--27; Phil 2:3--8; 1~Jn 3:16--18.
\item \work{Nova Vulgata Bibliorum Sacrorum editio}, the Gospel of Matthew as published by the Holy See at vatican.va: Mt 21:28--32.
\item \work{The Greek New Testament: SBL Edition}, ed. M.~W. Holmes, version~1.2, text and apparatus of Matthew at 21:29--31.
\end{itemize}

\subsection*{Fathers, Doctors and saints}
\begin{itemize}[itemsep=0pt,topsep=.25em]
\item St John Chrysostom, \work{Homilies on Matthew} 67.2--4, Nicene and Post-Nicene Fathers, 1st ser., vol.~10; \work{Homilies on Philippians} 5--7, Nicene and Post-Nicene Fathers, 1st ser., vol.~13 (both in the electronic text of the Christian Classics Ethereal Library).
\item St Jerome, \work{Commentarii in Matthaeum} III, on 21:27--32 (Migne, PL 26, Paris, 1845, cols.~155--156); \work{Commentarii in Hiezechielem} VI, on 18:21--32 (Migne, PL 25, Paris, 1845, cols.~179--181); \work{Commentaria in Danielem}, on 3:23--29 (PL 25, col.~509).
\item St Hilary of Poitiers, \work{Commentarius in Matthaeum} 21.11--15 (Migne, PL 9, Paris, 1844, cols.~1039--1041); \work{Tractatus super Psalmos}, ed. A. Zingerle, CSEL 22 (Vienna, 1891), on Ps~118, Zain 1--2 (pp.~418--419).
\item St Augustine, \work{Sermo} 82.14 and 87.10--11 (Sermons 32 and 37 on the New Testament Lessons), Nicene and Post-Nicene Fathers, 1st ser., vol.~6.
\item St Augustine, \work{Enarrationes in Psalmos} 17.27--28 and 24.1--9, Nicene and Post-Nicene Fathers, 1st ser., vol.~8; 118, s.~15.1--2 (Migne, PL 37, Paris, 1845, cols.~1540--1541).
\item St Augustine, \work{In Iohannis evangelium tractatus} 36.4, 48.4--6, 55.7, 84.1--2, and \work{In epistolam Iohannis ad Parthos tractatus} 5.11--12, Nicene and Post-Nicene Fathers, 1st ser., vol.~7.
\item St Gregory the Great, \work{Homiliae in Evangelia} 14.4, Latin (Wikisource text).
\item St Thomas Aquinas, \work{Super Evangelium Matthaei} c.~21 and \work{Super Evangelium Ioannis} c.~10, lect.~3 and 5, in \work{Opera}, editio altera Veneta, vol.~3 (Venice: Bettinelli, 1745), pp.~275--276, 660 and 668.
\item St Thomas Aquinas, \work{Super epistolam ad Philippenses} c.~2, lect.~1--3, in \work{In omnes S.~Pauli Apostoli epistolas commentaria}, vol.~2 (Li\`ege: Dessain, 1857), pp.~374--380.
\item St Thomas Aquinas, \work{Summa theologiae} I, q.~25, a.~3, obj.~3 and ad~3 (Leonine ed., vol.~4, Rome, 1888, pp.~293--294); II-II, q.~30, a.~4 (Corpus Thomisticum text); III, q.~84, a.~8, ad~1; q.~86, a.~1; q.~89, aa.~4--5 (English Dominican translation, Project Gutenberg text).
\item St Robert Bellarmine, \work{A Commentary on the Book of Psalms}, trans. J. O'Sullivan (1866, abridged), in the eCatholic2000 transcription: on Ps~24:4--9, Ps~118:49--50 and Ps~129:4--5.
\end{itemize}

\subsection*{Liturgical history and commentary}
\begin{itemize}[itemsep=0pt,topsep=.25em]
\item Blessed Ildefonso Schuster, \work{The Sacramentary (Liber Sacramentorum)}, vol.~3 (London: Burns Oates \& Washbourne, 1927), pp.~121--122 and 174--175.
\item H.~A. Wilson, ed., \work{The Gregorian Sacramentary under Charles the Great} (London, 1915), pp.~171--172; \work{The Gelasian Sacramentary} (Oxford, 1894); C.~L. Feltoe, ed., \work{Sacramentarium Leonianum} (Cambridge, 1896), for the Collect and the Prayer after Communion.
\end{itemize}

\subsection*{Magisterium}
\begin{itemize}[itemsep=0pt,topsep=.25em]
\item Second Vatican Council, Declaration on the Relation of the Church to Non-Christian Religions \work{Nostra aetate} (28~October 1965), no.~4, English text of the Holy See's website.
\end{itemize}

\subsection*{Chronology sources}
\begin{itemize}[itemsep=0pt,topsep=.25em]
\item \work{The Catholic Encyclopedia} (New York, 1908--1912): Reid, ``Captivities of the Israelites'' (vol.~3, 1908); Gigot, ``Book of Daniel'' (vol.~4, 1908); Schets, ``Ezekiel'' (vol.~5, 1909); Drum, ``Epistles of Saint John'' and Fonck, ``Gospel of St.~John'' (vol.~8, 1910); Jacquier, ``Gospel of St.~Matthew'' (vol.~10, 1911); Prat, ``St.~Paul'' (vol.~11, 1911); Vander Heeren, ``Epistle to the Philippians'' (vol.~12, 1911); Durand, ``The New Testament'' (vol.~14, 1912); read in the New Advent text.
\item United States Conference of Catholic Bishops, \work{New American Bible, Revised Edition}, introduction to the Psalms, \url{https://bible.usccb.org/bible/psalms/0}, read in its public delivery of 24~September 2026 (summarized; protected).
\item George Leo Haydock, note on Psalm~70:1, in the Douay--Rheims Bible with Haydock's commentary, reporting Usher's chronology; cited for the A.M.\ figure as the chronology corpus records it.
\end{itemize}
\par}
